\subsection{Reducción causal de la actualización bilateral}
\label{lib04:reduccion-causal}

La actualización unitaria anterior realiza el transporte del estado y de su
memoria en un mismo espacio. La observación de la primera componente admite
una descripción equivalente que conserva explícitamente la contribución del
registro. Esta descripción se obtiene del transporte ya construido, con los
mismos coeficientes y la misma orientación; la memoria eliminada de la lista
de variables reaparece como dependencia de los estados anteriores.

Fijemos los operadores \(A,D\) de \eqref{lib04:ad} y una realización
\(H\oplus H\). Para \(n\geq0\), sean
\begin{equation}
 v_{n+1}=Av_n-D^*m_n,\qquad
 m_{n+1}=Dv_n+A^*m_n.
 \label{lib04:recurrencia-memoria}
\end{equation}
La actualización conjunta es \(\mathscr R_a\) y conserva
\(\|v_n\|^2+\|m_n\|^2\).

\begin{proposition}[Eliminación exacta del registro auxiliar]
\label{lib04:prop-memoria-reducida}
Para todo \(n\geq0\), con la suma vacía interpretada como cero,
\begin{align}
 m_n&=(A^*)^n m_0+
       \sum_{k=0}^{n-1}(A^*)^{n-1-k}Dv_k,
       \label{lib04:memoria-explicita}\\
 v_{n+1}&=Av_n-D^*(A^*)^n m_0
       +\sum_{k=0}^{n-1}\mathscr K_{n-1-k}v_k,
       \label{lib04:evolucion-reducida}\\
 \mathscr K_j&=-D^*(A^*)^jD\quad(j\geq0).
       \label{lib04:nucleo-retardo}
\end{align}
La ecuación reducida, junto con \eqref{lib04:memoria-explicita}, es
equivalente a la recurrencia conjunta para los mismos datos \((v_0,m_0)\).
\end{proposition}
\begin{proof}
La primera identidad vale para \(n=0\). Si vale para \(n\), la segunda
ecuación de \eqref{lib04:recurrencia-memoria} proporciona
\[
 m_{n+1}=(A^*)^{n+1}m_0+
 \sum_{k=0}^{n-1}(A^*)^{n-k}Dv_k+Dv_n,
\]
que es la identidad en \(n+1\). Sustituirla, en profundidad \(n\), en
la primera ecuación demuestra \eqref{lib04:evolucion-reducida}.
Recíprocamente, si \(v_n\) satisface esta última y se define \(m_n\)
por \eqref{lib04:memoria-explicita}, la misma diferencia de sumas prueba
\(m_{n+1}=Dv_n+A^*m_n\); la otra ecuación se recupera por sustitución.
\end{proof}

El núcleo \(\mathscr K_j\) describe una transferencia hacia el registro,
\(j\) actualizaciones dentro de él y una reentrada en la componente leída.
El término \(-D^*(A^*)^nm_0\) conserva la memoria inicial. Su supresión
corresponde al caso particular \(m_0=0\), no a una identidad del transporte
general. En \eqref{lib04:evolucion-reducida} intervienen únicamente el
estado actual y los estados anteriores. La causalidad de esta reducción
procede del orden de la recurrencia conjunta.

La normalización del balance proporciona además una estimación explícita.
Como \(Q_-=aI-U\), se tiene
\(\|A\|^2\leq r_a=(a+1)^3/((a+1)^3+a^3)<1\), mientras que
\(D^*D\leq I\). Por tanto,
\begin{equation}
 \|\mathscr K_j\|\leq r_a^{j/2},\qquad
 \sum_{j\geq0}\|\mathscr K_j\|
 \leq\frac{1}{1-\sqrt{r_a}}.
 \label{lib04:cota-nucleo-retardo}
\end{equation}
Esta cota controla el peso de la memoria en la ecuación reducida; la norma
del estado conjunto sigue conservándose por la actualización unitaria.

La proposición es la realización temporal de la reducción de Schur de
\cref{mem:prop-schur}. En las series formales
\(V(z)=\sum_{n\geq0}v_nz^n\) y
\(M(z)=\sum_{n\geq0}m_nz^n\), las recurrencias dan
\(V=v_0+zAV-zD^*M\) y \(M=m_0+zDV+zA^*M\). Eliminando \(M\),
\begin{equation}
 V(z)=\bigl[I-zA+z^2D^*(I-zA^*)^{-1}D\bigr]^{-1}
       \bigl[v_0-zD^*(I-zA^*)^{-1}m_0\bigr].
 \label{lib04:schur-bilateral}
\end{equation}
Las inversas formales existen porque sus términos constantes son identidades;
las series de los estados y el resolvente conjunto también convergen para
\(|z|<1\). La fuente inicial y el operador efectivo se transforman juntos.
Un lector que actúe también sobre \(m_n\) conserva su contribución mediante
\eqref{lib04:memoria-explicita}, conforme a \eqref{mem:schur-lector}.

La reentrada aquí descrita actualiza un registro entrante. La política de
archivo que sigue conserva, en cambio, cada complemento en una coordenada
distinta. Ambas operaciones poseen su propia ley de conservación. El núcleo
\(\mathscr K_j\), los propagadores de Green y una fuerza dimensional son
objetos de tipos diferentes: el primero elimina una componente de estado,
los segundos resuelven ecuaciones de fuente con condiciones de soporte, y
la fuerza añade una realización dinámica y sus unidades.
